% engines: lualatex
% expect-warnings: 0
% expect-errors: 0
\documentclass{article}
\usepackage[margin=1cm]{geometry}
\usepackage[kernel = lua]{../../latex/pgf-go}
\newcommand\case[1]{\par\medskip\noindent\textbf{#1}\par\nopagebreak}
\pagestyle{empty}
\begin{document}

% game.sgf: 9x9; B E5, W G7, B pass, W C3, B G5.

\case{pgfgolayout, columns = 3: boards of different height, valign = top (initially)}
\begin{pgfgolayout}[columns = 3]
	\begin{goban}[5, scale = 0.5]\stones[b]{C3}\end{goban}
	\begin{goban}[5x3, scale = 0.5]\stones[w]{C2}\end{goban}
	\begin{goban}[9, scale = 0.5]\stones[b]{E5}\end{goban}
\end{pgfgolayout}

\case{the same with valign = bottom: the boards end at the same height}
\begin{pgfgolayout}[columns = 3, valign = bottom]
	\begin{goban}[5, scale = 0.5]\stones[b]{C3}\end{goban}
	\begin{goban}[5x3, scale = 0.5]\stones[w]{C2}\end{goban}
	\begin{goban}[9, scale = 0.5]\stones[b]{E5}\end{goban}
\end{pgfgolayout}

\case{sgf diagrams, partial = auto, columns = 3, valign = bottom: the captions start at the same height; captions in several lines}
\sgfdiagrams[step = 2, goban = {partial = auto, scale = 0.5},
	caption = {font = \large}, layout = {columns = 3, width = 8cm, valign = bottom}]{game.sgf}

\case{the same with width = natural}
\sgfdiagrams[step = 2, goban = {partial = auto, scale = 0.5},
	layout = {columns = 3, width = natural, valign = bottom}]{game.sgf}

\case{sgf problem, statement = false, columns = 3, fit, valign = bottom}
\sgfproblem[statement = false, goban = {partial = auto}, layout = {columns = 3, fit, valign = bottom}]{problem.sgf}

\end{document}
